% =============================================================================
% 附录——编辑说明 — Paratextual frame and reconstruction methodology
% =============================================================================
\cleardoublepage
\thispagestyle{empty}
\chapter*{附录——编辑说明}
\addcontentsline{toc}{chapter}{附录——编辑说明}
\markboth{编辑说明}{}

\begingroup

本版依据患者1,260原始笔记本的法证誊本编订，该笔记本于1979年由验尸机构保管，附于巴巴塞纳收容院病历。

\vspace{0.6em}

\noindent\textbf{方法:}
采用的编辑标准是对原始文件的最大忠实，保留其碎片化结构、空缺、重复和内部矛盾。这里汇集的78个碎片遵循大致的写作顺序，通过载体、墨迹和笔迹的相对年代测定重建，依据所附法证报告（DOC-03）。

\vspace{0.6em}

\noindent\textbf{副文本装置:}
\begin{itemize}
  \item \textbf{[编者按]} —— 编辑对内容、来源或重建限度的评论。
  \item \textbf{[无法辨认]} —— 原文因载体退化而无法破译的段落。
  \item \textbf{[缺文]} —— 原始笔记本中缺失的书页或段落。
  \item \textbf{[原文如此]} —— 对原文错误或特质的逐字转录。
\end{itemize}

\vspace{0.6em}

\noindent\textbf{限度:}
本版不声称完整。原始日记约有847页；读者拿到的是经过编选的选本，按五个主题维度（MEM、DOC、LUC、CONT）加尾声组织。因冗余、不可逆退化或保护第三方隐私而被删除的碎片，已在档案清单中按法证封印登记。

\vspace{0.6em}

\noindent\textbf{最后警告:}
本书是重建，不是摹本复制。编辑选择——排序、筛选、注释——反映了出版日期的研究状态。未来的研究者若获得额外的非破坏性技术，可以扩展或修正这里呈现的内容。

\vspace{1.2em}

\begin{flushright}
--- 编委会\par
OpenCode Books，克拉特乌斯\par
2026年7月
\end{flushright}

\endgroup
\clearpage
